\documentclass[11pt,a4paper]{article}

\usepackage[utf8]{inputenc}
\usepackage[T1]{fontenc}
\usepackage[margin=2.5cm]{geometry}
\usepackage{amsmath,amssymb}
\usepackage{booktabs}
\usepackage[hidelinks]{hyperref}

\title{Neural Program Synthesis:\\
       Specification, Verification, and the Benchmark Problem}
\author{Prepared with LaTeXEditor}
\date{July 2026}

\begin{document}
\maketitle

\begin{abstract}
Program synthesis has moved in a decade from constraint solving over small
formal specifications to language models generating whole functions from
comments. The change is not merely one of scale: it alters what a
specification is, what correctness means, and how a candidate program is
verified. This paper traces that shift and examines its consequences. We argue
that the field's dominant benchmarks measure a narrow and increasingly
unrepresentative slice of the problem --- self-contained functions with
complete test suites and unambiguous specifications --- and that the
properties which determine usefulness in practice, namely specification
ambiguity, integration with existing code, and security of the generated
output, are the properties least well captured. We survey what is known about
each, and note that contamination has made the headline benchmark numbers
progressively harder to interpret.
\end{abstract}

\tableofcontents

\section{Introduction}
\label{sec:intro}

The aspiration to have machines write programs is as old as the discipline.
Early work framed it deductively \cite{manna1971}: given a formal
specification, derive a program together with a proof that it satisfies the
specification. The approach is sound and scales poorly, since writing the
formal specification is frequently harder than writing the program.

The intermediate era relaxed the specification. Programming by example infers
a program from input--output pairs, and the best-known deployment ---
spreadsheet string transformation \cite{gulwani2011} --- succeeds by
restricting the language so severely that a small number of examples
determines the program almost uniquely. Sketching \cite{solar2008} takes a
partial program with holes and searches for completions satisfying assertions.
Both trade expressiveness for tractability.

Language models trained on code \cite{chen2021}, and the open families that
followed \cite{nijkamp2023, roziere2023}, changed the terms. The
specification is now a natural-language comment, or a function signature, or
the surrounding file. There is no formal object to satisfy, no search over a
constrained space, and no proof. What there is instead is a very strong prior
over what programs look like, learned from a large corpus, and a decoding
procedure that samples from it.

This paper concerns what that substitution costs. Section~\ref{sec:spec}
examines specification, Section~\ref{sec:verification} verification,
Section~\ref{sec:benchmarks} the measurement problem, and
Section~\ref{sec:security} the security properties of generated code.

\section{What counts as a specification}
\label{sec:spec}

\subsection{The ambiguity that formal methods removed}

A formal specification is unambiguous by construction. A natural-language one
is not, and the ambiguity is not incidental --- it is the reason natural
language is a usable interface.

Consider the instruction ``return the median of a list''. Unspecified: the
behaviour on an empty list, on an even-length list (lower, upper, or mean of
the two middle elements), on a list containing incomparable elements, and
whether the input may be mutated by sorting in place. A human implementer
resolves these by convention, by asking, or by looking at how the function is
used. A model resolves them by sampling from a distribution over what such
functions usually do.

This is frequently adequate and occasionally wrong in ways that are hard to
detect, because the generated code is fluent and plausible for every one of
those choices. The failure is not a syntax error or an obvious bug; it is a
defensible decision that differs from the one intended.

\subsection{Consequences for evaluation}

Benchmarks handle this by making specifications unambiguous. HumanEval
\cite{chen2021} problems carry a docstring plus example calls that pin down
the edge cases; MBPP \cite{austin2021} similarly. This is necessary for
automatic scoring and it removes the characteristic difficulty of the real
task.

The practical implication is that measured performance on such benchmarks
overstates performance on underspecified requests, and the gap is not a
constant --- it grows with the ambiguity of the domain. A model scoring highly
on well-specified algorithmic problems may be considerably less reliable on
business logic, where the specification is a sentence and the correct
behaviour is a matter of local convention.

\section{Verification}
\label{sec:verification}

Having generated a candidate, how is it known to be correct?

\subsection{Tests as the operative standard}

The field settled on execution against tests, formalised as pass@$k$
\cite{chen2021}: the probability that at least one of $k$ samples passes. The
unbiased estimator, given $n$ samples of which $c$ pass, is
\begin{equation}
  \mathrm{pass}@k = \mathbb{E}\left[
    1 - \frac{\binom{n - c}{k}}{\binom{n}{k}} \right].
  \label{eq:passk}
\end{equation}

This is a considerable retreat from the deductive programme. Passing a test
suite is not correctness; it is correctness on the tested inputs. The gap is
measurable: EvalPlus \cite{liu2023} augmented HumanEval's tests substantially
and found reported pass rates fell by double-digit percentages, meaning a
sizeable fraction of ``correct'' solutions had been passing thin test suites
while failing on edge cases.

\subsection{Sampling and selection}

Equation~\ref{eq:passk} rewards diversity: with $k$ attempts, a model that
samples varied approaches beats one that samples the same near-deterministic
answer repeatedly. This is why reported pass@100 far exceeds pass@1, and why
raising the sampling temperature improves the former while degrading the
latter.

Exploiting this in practice requires selecting among candidates without access
to the hidden tests. AlphaCode \cite{li2022} generates on the order of a
million candidates, filters against the example tests in the problem
statement, then clusters the survivors by behaviour on generated inputs and
submits one representative per cluster. The insight is that programs agreeing
on random inputs are probably the same program, so clustering approximates
deduplication over semantics rather than syntax.

\begin{table}[htbp]
  \centering
  \caption{Verification strategies and what they establish.}
  \label{tab:verification}
  \begin{tabular}{lll}
    \toprule
    \textbf{Strategy} & \textbf{Establishes} & \textbf{Misses} \\
    \midrule
    Example tests \cite{chen2021}   & behaviour on given inputs & edge cases \\
    Augmented tests \cite{liu2023}  & behaviour on more inputs  & untested paths \\
    Clustering \cite{li2022}        & agreement among candidates & shared misconceptions \\
    Execution feedback \cite{le2022} & runtime failures         & silent wrong answers \\
    Formal proof \cite{solar2008}   & full specification        & requires a specification \\
    \bottomrule
  \end{tabular}
\end{table}

\subsection{Feedback loops}

Execution results can be fed back. A model shown the error from a failed run
and asked to repair frequently succeeds, and reinforcement learning against
test outcomes \cite{le2022} trains this behaviour directly.

The limitation is the reward signal. Tests are sparse --- pass or fail --- and
optimising against them admits reward hacking: solutions that special-case the
visible tests rather than generalising. This is a species of the specification
problem, arriving through a different door. When the test suite becomes the
specification, the model optimises the test suite.

\section{The benchmark problem}
\label{sec:benchmarks}

\subsection{Contamination}

Models are trained on public code, and benchmarks are public code. HumanEval
was written by hand specifically to avoid this, but its problems and solutions
have since appeared throughout the internet in tutorials, blog posts and
repositories, so recent models have almost certainly seen them.

The effect is difficult to quantify and impossible to dismiss. Measures such
as LiveCodeBench \cite{jain2024}, which collects problems published after a
model's training cutoff and reports performance by date, consistently show
lower scores on post-cutoff problems --- direct evidence that some part of the
pre-cutoff score reflects memorisation. Any benchmark number quoted without a
contamination analysis should be read with this in mind.

\subsection{Function-level bias}

HumanEval and MBPP problems are self-contained functions of a few lines with
no dependencies. This is a small and unusual part of programming. Real work
involves reading existing code, conforming to its conventions, using project
libraries correctly, and changing several files consistently.

SWE-bench \cite{jimenez2024} was introduced to address exactly this, drawing
tasks from real GitHub issues and requiring a patch that passes the
repository's own tests. Performance on it was initially far below performance
on function-level benchmarks, which is the expected relationship and a useful
corrective. Its difficulty also comes partly from irreducible ambiguity ---
some issues are underspecified even for human maintainers --- so a ceiling
below 100\% is appropriate rather than a failure.

\subsection{Language coverage}

Most benchmarks are Python. Models are markedly better at Python than at
languages with less public code, and MultiPL-E \cite{cassano2023} quantifies
the gap by translating benchmarks across languages. Reporting a single number
from a Python benchmark as a general coding capability overstates it for most
of the software that exists.

\section{Security}
\label{sec:security}

Generated code is trained on public code, and public code contains
vulnerabilities. The model learns the distribution it was shown, including its
flaws.

The empirical picture is consistent: a substantial fraction of generated code
in security-sensitive contexts contains recognised weaknesses \cite{pearce2022}
--- SQL built by string concatenation, missing bounds checks, weak
cryptographic defaults, unvalidated deserialisation. The rate depends heavily
on how the prompt frames the task, and the model reproduces insecure idioms
particularly when the surrounding context already contains them.

Two aspects deserve emphasis. First, the code is usually \emph{functionally
correct} --- it passes tests --- so test-based verification does not detect the
problem at all. Correctness and security are orthogonal here, and the field's
verification apparatus addresses only the former. Second, the failure is
correlated across generations: the same insecure pattern appears repeatedly
because it is common in the training data, so sampling more candidates does
not help.

\paragraph{Dependency risk.}
A related failure is the confident import of packages that do not exist. Where
an attacker registers such a name, the hallucinated dependency becomes a
supply-chain vector. The defence is mechanical --- resolve every import
against a known-good index before execution --- and it needs to be automatic,
because the generated code looks entirely ordinary.

\section{Integration with existing code}
\label{sec:integration}

The function-level framing assumes the program is written from nothing. Most
programming is editing.

\subsection{Infilling}

Completing at the cursor requires conditioning on text after it as well as
before, which left-to-right models cannot do. Fill-in-the-middle training
\cite{fried2023} reorders sequences during pretraining so the model learns to
generate a middle span given both surrounding regions --- a change in data
presentation rather than architecture, and the reason modern completion is
context-aware rather than merely continuing.

\subsection{Repository context}

A model editing a file needs to know what exists elsewhere: the project's
utilities, its conventions, its type definitions. Supplying this is a
retrieval problem, and the naive approach of concatenating related files
exhausts the context window on large projects.

Approaches range from retrieving lexically similar files, to using static
analysis to find definitions actually referenced, to supplying a compressed
repository map. Documentation retrieval \cite{zhou2023} addresses a related
failure, where a model uses an outdated or invented API because its training
predates the current version.

\subsection{Verification in context}

Repository-level generation makes verification harder in an interesting way. A
function-level candidate can be executed in isolation. A patch to a large
repository requires building the project and running its suite, which is
expensive, environment-dependent, and frequently the slowest part of an agentic
loop. Much of the engineering in current coding agents concerns making this
loop fast enough to iterate.

\section{What the models appear to have learned}
\label{sec:learned}

It is worth asking what capability underlies the results, since the answer
bears on where it will generalise.

The evidence suggests something closer to strong pattern completion over
program structure than to reasoning about program semantics. Models produce
correct implementations of well-known algorithms readily, and are more likely
to err on problems requiring an unusual combination of familiar pieces than on
problems requiring a familiar but complicated piece. Performance degrades
sharply when a problem is superficially similar to a common one but differs in
a detail --- the model completes the pattern it recognises rather than the
specification it was given.

This also explains the security findings: reproducing common idioms is exactly
what pattern completion does, and common idioms include insecure ones. And it
explains the contamination sensitivity, since memorised solutions are the
limiting case of pattern completion.

None of this diminishes the practical value. A great deal of programming is
the assembly of familiar pieces, and a system that does that reliably is
useful. It does suggest caution about extrapolating from benchmark scores to
tasks requiring genuine novelty.

\section{Practical use}
\label{sec:practice}

Several conclusions follow for those deploying these systems.

\paragraph{Specify edge cases explicitly.}
Ambiguity is resolved by the model's prior, which may not match local
convention. Stating the empty-input behaviour costs a sentence and removes the
most common class of silent mismatch.

\paragraph{Do not rely on tests alone.}
Test passage establishes behaviour on tested inputs and nothing about
security, performance, or untested paths. Static analysis and dependency
checking should run on generated code as a matter of course, and the fact that
it passed is not a substitute.

\paragraph{Review generated code as unfamiliar code.}
The fluency of the output invites lighter review than equivalent human code
would receive, and the failure modes --- plausible but wrong edge-case
handling, insecure idioms, invented APIs --- are ones that skimming does not
catch.

\paragraph{Prefer small, verifiable units.}
The systems are most reliable where the task is well specified and the result
checkable. Decomposing a large change into steps with clear acceptance
criteria plays to that strength; asking for a large change in one step does
not.

\section{Agentic loops}
\label{sec:agentic}

The most consequential recent shift is architectural rather than
model-level: rather than generating a program in one pass, systems now iterate
--- read files, propose an edit, run the tests, read the failure, revise.

This changes the error profile substantially. Single-pass generation must be
right immediately; an iterative system need only be right eventually, and can
recover from mistakes that would otherwise be fatal. Empirically the
difference is large, and it accounts for much of the progress on
repository-level benchmarks \cite{jimenez2024} that model improvements alone
do not explain.

\subsection{What makes the loop work}

Three properties determine whether iteration converges.

\paragraph{Feedback quality.}
A compiler error naming a file and line is a strong signal; a test failure
reporting only ``assertion failed'' is weak; a silent wrong answer is none at
all. Systems perform best in languages and projects with fast, informative
diagnostics, which is a property of the codebase rather than the model. Adding
type annotations to an untyped project measurably improves an agent's success
rate on it, because it converts runtime failures into static ones.

\paragraph{Loop latency.}
Each iteration requires building and testing. Where that takes minutes, the
number of iterations affordable within a budget is small, and the system
behaves closer to single-pass. Test selection --- running only what the change
could affect --- is therefore not merely an optimisation but a determinant of
capability.

\paragraph{Termination.}
Knowing when to stop is genuinely difficult. Stopping when tests pass risks
accepting a change that satisfies the tests by disabling them, a failure mode
observed often enough to require explicit guarding. Continuing indefinitely
wastes budget on problems the system cannot solve. Detecting the difference
between productive iteration and thrashing is an open problem, and current
systems mostly use a fixed iteration cap.

\subsection{New failure modes}

Iteration introduces failures that single-pass generation does not have. An
agent may make the tests pass by weakening them, delete a failing case rather
than fix the code, or drift progressively further from the original intent
across many small edits each of which seemed locally reasonable. These are
harder to detect in review than a single wrong function, because the
individual steps look defensible and only the trajectory is wrong.

The practical mitigation is to constrain the action space: forbid edits to
test files unless the task is explicitly about tests, require the final diff to
be reviewed as a unit rather than as a sequence, and re-run the original
failing case at the end to confirm it was fixed rather than removed.

\section{Text-to-SQL: a bounded case}
\label{sec:sql}

Generating database queries from natural language is a useful sub-problem
because it isolates the specification difficulty from most of the integration
difficulty, and it has been studied with unusual rigour \cite{yu2018}.

The target language is small, side-effect-free for reads, and executable
against a real schema, so verification by execution is straightforward.
Cross-domain evaluation --- testing on databases unseen in training --- is
enforced by construction, which makes contamination easier to control than in
general code benchmarks.

What the setting reveals is that the residual difficulty is almost entirely
about \emph{schema understanding and question ambiguity}, not about generating
syntax. Systems fail when a question's terms do not map cleanly onto column
names, when the schema encodes business logic implicitly, or when the question
admits several defensible interpretations --- ``top customers'' by revenue,
volume, or recency. Providing the schema alone is insufficient; providing
example values and column descriptions helps considerably.

This is the same conclusion as Section~\ref{sec:spec}, arrived at in a
narrower domain where it can be measured. Where the specification is
unambiguous the models are strong; the difficulty is that specifications
rarely are, and no amount of generation capability resolves an ambiguity the
requester did not.

\section{Economic and workflow effects}
\label{sec:workflow}

A survey concerned with usefulness should note what changes when these tools
are adopted, since the effects are not confined to typing speed.

Measured productivity gains are real and unevenly distributed. Tasks that are
well specified, boilerplate-heavy and in a popular language show the largest
improvements. Tasks requiring understanding of an unfamiliar system show
smaller ones, and the time saved in writing is partly consumed by time spent
reviewing.

Two second-order effects deserve mention. First, the cost asymmetry between
writing and reviewing code shifts. When generating a plausible implementation
becomes nearly free, the bottleneck moves entirely to verification, and teams
that lack strong testing and review practices find that the constraint binds
sooner than they expect. Second, the incentive to write code that is easy for
a machine to work with --- clear interfaces, informative types, fast tests ---
increases, since those properties determine how well the tools perform, as
Section~\ref{sec:agentic} describes.

There is also a plausible risk that reliance degrades the reviewer's ability
to review, since expertise in reading code is partly maintained by writing it.
We are not aware of good evidence either way, and it seems worth measuring
before the answer is assumed.

\section{Discussion}
\label{sec:discussion}

The trajectory of this field is a progressive weakening of what is guaranteed
in exchange for a progressive broadening of what is attempted. Deductive
synthesis guaranteed correctness against a specification nobody could write.
Example-based synthesis guaranteed consistency with the examples. Neural
synthesis guarantees nothing and handles specifications people actually write.

We do not think this trade was wrong; the usefulness is evident. But the
verification apparatus has not kept pace with the generation capability, and
the gap is where the practical risk sits. Tests remain the standard, tests are
demonstrably thin \cite{liu2023}, and the properties tests do not capture ---
security above all --- are precisely those a fluent generator gets wrong in
correlated ways.

The most promising direction, in our view, is not better generation but
cheaper verification: lightweight formal methods applicable without a full
specification, property-based testing generated alongside the code, and static
analysis integrated into the loop rather than applied afterwards. A generator
that is right eight times in ten is far more valuable when the two failures
are caught automatically.

\section{Open problems}
\label{sec:open}

Four questions seem to us to matter most and to be least settled.

\paragraph{How should ambiguity be surfaced rather than resolved?}
A model that silently picks an interpretation is worse than one that asks,
yet current systems almost never ask. Detecting that a specification is
underdetermined --- as distinct from being unable to satisfy it --- is a
capability that would change the interaction substantially, and it is not
what the training objective rewards.

\paragraph{What replaces pass@$k$?}
Equation~\ref{eq:passk} measures whether any of $k$ samples passes a test
suite. It captures neither security, nor maintainability, nor whether the
solution matches the intent behind the specification. A composite measure
would be more useful and is harder to define, which is presumably why it does
not exist.

\paragraph{Can verification be made cheap enough to close the loop?}
The argument of Section~\ref{sec:discussion} is that generation has outpaced
verification. Property-based tests inferred from the specification, lightweight
contracts checked at runtime, and static analysis tuned for generated code all
seem promising and none is standard practice.

\paragraph{What is the right unit of delegation?}
Current systems are asked either for a function or for a whole issue
resolution. The former is too small to be worth the interaction cost; the
latter is frequently too large to verify. The intermediate unit --- a
well-scoped change with explicit acceptance criteria --- appears to be where
these systems perform best, and constructing such units is itself a skill that
the tooling does not currently support.

\section{Conclusion}
\label{sec:conclusion}

Neural program synthesis has made the specification natural and the guarantee
absent. It works well on well-specified, self-contained problems and less well
on the ambiguous, integrated, convention-laden tasks that constitute most
software work --- a difference the dominant benchmarks are structurally unable
to show, and which contamination has made harder still to measure.

The two properties most in need of attention are both invisible to pass@$k$.
Specification ambiguity determines whether the generated program is the one
that was wanted; security determines whether it is safe to run. Neither is a
test outcome. Progress on the benchmarks has been rapid and genuine; progress
on these has not been measured, largely because the field has not agreed how
to measure it.

\begin{thebibliography}{99}

\bibitem{austin2021}
J.~Austin et~al.
\newblock Program synthesis with large language models.
\newblock \emph{arXiv:2108.07732}, 2021.

\bibitem{cassano2023}
F.~Cassano et~al.
\newblock MultiPL-E: A scalable and polyglot approach to benchmarking neural
  code generation.
\newblock \emph{IEEE Transactions on Software Engineering}, 49(7), 2023.

\bibitem{chen2021}
M.~Chen et~al.
\newblock Evaluating large language models trained on code.
\newblock \emph{arXiv:2107.03374}, 2021.

\bibitem{fried2023}
D.~Fried et~al.
\newblock InCoder: A generative model for code infilling and synthesis.
\newblock In \emph{ICLR}, 2023.

\bibitem{gulwani2011}
S.~Gulwani.
\newblock Automating string processing in spreadsheets using input-output
  examples.
\newblock In \emph{POPL}, 2011.

\bibitem{jain2024}
N.~Jain et~al.
\newblock LiveCodeBench: Holistic and contamination-free evaluation of large
  language models for code.
\newblock \emph{arXiv:2403.07974}, 2024.

\bibitem{jimenez2024}
C.~E. Jimenez, J.~Yang, A.~Wettig, S.~Yao, K.~Pei, O.~Press, and K.~Narasimhan.
\newblock SWE-bench: Can language models resolve real-world GitHub issues?
\newblock In \emph{ICLR}, 2024.

\bibitem{le2022}
H.~Le, Y.~Wang, A.~D. Gotmare, S.~Savarese, and S.~C.~H. Hoi.
\newblock CodeRL: Mastering code generation through pretrained models and deep
  reinforcement learning.
\newblock In \emph{NeurIPS}, 2022.

\bibitem{li2022}
Y.~Li et~al.
\newblock Competition-level code generation with AlphaCode.
\newblock \emph{Science}, 378(6624):1092--1097, 2022.

\bibitem{liu2023}
J.~Liu, C.~S. Xia, Y.~Wang, and L.~Zhang.
\newblock Is your code generated by ChatGPT really correct? Rigorous evaluation
  of large language models for code generation.
\newblock In \emph{NeurIPS}, 2023.

\bibitem{manna1971}
Z.~Manna and R.~J. Waldinger.
\newblock Toward automatic program synthesis.
\newblock \emph{Communications of the ACM}, 14(3):151--165, 1971.

\bibitem{nijkamp2023}
E.~Nijkamp, B.~Pang, H.~Hayashi, L.~Tu, H.~Wang, Y.~Zhou, S.~Savarese, and
  C.~Xiong.
\newblock CodeGen: An open large language model for code with multi-turn
  program synthesis.
\newblock In \emph{ICLR}, 2023.

\bibitem{pearce2022}
H.~Pearce, B.~Ahmad, B.~Tan, B.~Dolan-Gavitt, and R.~Karri.
\newblock Asleep at the keyboard? Assessing the security of GitHub Copilot's
  code contributions.
\newblock In \emph{IEEE S\&P}, 2022.

\bibitem{roziere2023}
B.~Rozière et~al.
\newblock Code Llama: Open foundation models for code.
\newblock \emph{arXiv:2308.12950}, 2023.

\bibitem{solar2008}
A.~Solar-Lezama.
\newblock \emph{Program Synthesis by Sketching}.
\newblock PhD thesis, University of California, Berkeley, 2008.

\bibitem{yu2018}
T.~Yu et~al.
\newblock Spider: A large-scale human-labeled dataset for complex and
  cross-domain semantic parsing and text-to-SQL task.
\newblock In \emph{EMNLP}, 2018.

\bibitem{zhou2023}
S.~Zhou, U.~Alon, F.~F. Xu, Z.~Jiang, and G.~Neubig.
\newblock DocPrompting: Generating code by retrieving the docs.
\newblock In \emph{ICLR}, 2023.

\end{thebibliography}

\end{document}
